\begin{thebibliography}{9}
\bibitem{Euler}
Euler, L. (1750/2007). \emph{A double demonstration of a theorem of Newton, which
gives a relation between the coefficient of an algebraic equation and the sums
of the powers of its roots} (J. Bell, Trans.). Original: \emph{Opuscula varii
argumenti}, 2, 108--120, E153. arXiv:0707.0699. \S\S3, 5--7 are used only for
the classical Newton identities, which are also derived in the proof.
\url{https://arxiv.org/abs/0707.0699}
\bibitem{ET}
Elsholtz, C., \& Tao, T. (2013). Counting the number of solutions to the
Erd\H{o}s--Straus equation on unit fractions. \emph{Journal of the Australian
Mathematical Society, 94}, 50--105. arXiv:1107.1010, \S2, Propositions 2.2 and
2.6. The Type-I/II identities are inherited; no averaged theorem is used as a
pointwise occupancy statement. \url{https://arxiv.org/abs/1107.1010}
\bibitem{Milne}
Milne, J. S. (2020). \emph{Algebraic number theory}, version 3.08.
Chapter 7, Newton's lemma, supplies the standard local lifting background.
The exact quadratic lift used here is given in the proof.
\url{https://www.jmilne.org/math/CourseNotes/ANT.pdf}
\bibitem{Previous}
The Clankers. (2026). \emph{Original integral completion of the literal
ES--Fabel inverse}. Preserved Turn 7 source, 20 September 2026,
\texttt{core.tex} in the supplied Turn 7 integral-completion directory.
The original IC4--IC7 quartic coordinates, IC14 local residue calculation, and
the foreign-prime source-domain obstruction are used at their stated scope.
The exact source hashes and retrieval locators are recorded in
\texttt{source\_reading.json}. Earlier work remains separately preserved in the
cumulative archive; it is not rewritten as part of this continuation.
\end{thebibliography}
